\documentclass[12pt,a4paper]{article}
\usepackage[utf8]{inputenc}
\usepackage[margin=1in]{geometry}
\usepackage{amsmath}
\usepackage{amsfonts}
\usepackage{amssymb}
\usepackage{graphicx}
\usepackage{listings}
\usepackage{xcolor}
\usepackage{tikz}
\usepackage{algorithm}
\usepackage{algorithmic}
\usepackage{fancyhdr}
\usepackage{hyperref}
\usepackage{enumitem}

% Code styling
\lstset{
    backgroundcolor=\color{gray!10},
    basicstyle=\ttfamily\footnotesize,
    breakatwhitespace=false,
    breaklines=true,
    commentstyle=\color{green!60!black},
    keywordstyle=\color{blue},
    stringstyle=\color{red},
    numberstyle=\tiny\color{gray},
    numbers=left,
    stepnumber=1,
    tabsize=2,
    frame=single,
    rulecolor=\color{black!30}
}

% Header and footer
\pagestyle{fancy}
\fancyhf{}
\fancyhead[L]{Complete Algorithm Logic Flow}
\fancyhead[R]{3D Boolean Operations}
\fancyfoot[C]{\thepage}

\title{Complete Algorithm Logic Flow: \\
\large Sophisticated 3D Depth-Based Polygon Classification \\
with Farthest Point Detection and Face Association}
\author{3D Boolean Operations System}
\date{\today}

\begin{document}

\maketitle

\tableofcontents
\newpage

\section{Algorithm Overview}

This document describes a sophisticated algorithm for depth-based polygon classification using 3D geometric analysis. The algorithm processes 2D polygon projections from 3D faces and uses multi-point sampling to determine which polygons are "farther" along a projection normal direction, applying boolean operations accordingly.

\subsection{Key Features}
\begin{itemize}
    \item \textbf{Multi-point sampling}: Uses 3×3 grid sampling for accurate depth analysis
    \item \textbf{Farthest point detection}: Finds the point that extends farthest along projection normal
    \item \textbf{Face association}: Associates intersection regions with their originating faces
    \item \textbf{Separate intersection tracking}: Creates distinct intersection rectangles for each polygon pair
\end{itemize}

\section{Algorithm Data Structures}

\subsection{Polygon Data Structure}
Each polygon is stored as a dictionary containing:
\begin{lstlisting}[language=Python]
polygon_data = {
    'polygon': shapely.Polygon,      # The 2D polygon geometry
    'name': str,                     # Face identifier (e.g., "Face_4")
    'normal': np.array,              # 3D normal vector
    'parent_face': OpenCASCADE_Face, # 3D face reference
    'original_index': int            # Original position index
}
\end{lstlisting}

\subsection{Array Definitions}
\begin{itemize}
    \item \textbf{array\_A}: Initial collection of all polygons to be processed
    \item \textbf{array\_B}: Unclassified polygons (no intersections found yet)
    \item \textbf{array\_C}: Classified results (intersections + remaining polygon pieces)
\end{itemize}

\section{Complete Algorithm Logic Flow}

\subsection{Step 1: Array\_A Population}

\begin{algorithm}
\caption{Array\_A Population}
\begin{algorithmic}[1]
\STATE Initialize $\text{array\_A} \leftarrow \emptyset$
\FOR{each $(polygon_i, name_i, normal_i)$ in $polygons_{2D}$}
    \STATE $parent\_face_i \leftarrow face\_associations[i]$
    \STATE $polygon\_data_i \leftarrow \{$
    \STATE \quad $'polygon': polygon_i,$
    \STATE \quad $'name': name_i,$
    \STATE \quad $'normal': normal_i,$
    \STATE \quad $'parent\_face': parent\_face_i,$
    \STATE \quad $'original\_index': i$
    \STATE $\}$
    \STATE $\text{array\_A} \leftarrow \text{array\_A} \cup \{polygon\_data_i\}$
\ENDFOR
\end{algorithmic}
\end{algorithm}

\textbf{Purpose}: Load all 2D polygons into array\_A with complete metadata including parent face associations for 3D depth calculations.

\subsection{Step 2: Array Initialization and Seeding}

\begin{algorithm}
\caption{Array Initialization}
\begin{algorithmic}[1]
\STATE Initialize $\text{array\_B} \leftarrow \emptyset$
\STATE Initialize $\text{array\_C} \leftarrow \emptyset$
\IF{$\text{array\_A} \neq \emptyset$}
    \STATE $first\_polygon \leftarrow \text{array\_A.pop}(0)$
    \STATE $\text{array\_B} \leftarrow \{first\_polygon\}$
\ENDIF
\end{algorithmic}
\end{algorithm}

\textbf{Purpose}: Initialize empty result arrays and seed array\_B with the first polygon to start the iterative comparison process.

\subsection{Step 3: Main Processing Loop}

\subsubsection{Step 3a: Polygon Selection}

\begin{algorithm}
\caption{Main Processing Loop}
\begin{algorithmic}[1]
\WHILE{$\text{array\_A} \neq \emptyset$}
    \STATE $P_i\_data \leftarrow \text{array\_A.pop}(0)$
    \STATE $P_i \leftarrow P_i\_data['polygon']$
    \STATE $P_i\_name \leftarrow P_i\_data['name']$
    \STATE $P_i\_parent\_face \leftarrow P_i\_data['parent\_face']$
    \STATE $intersection\_found \leftarrow \text{False}$
\end{algorithmic}
\end{algorithm}

\subsubsection{Step 3b: Intersection Testing}

\begin{algorithm}
\caption{Intersection Testing Against array\_B}
\begin{algorithmic}[1]
\FOR{each $P_j\_data$ in $\text{array\_B}$}
    \STATE $P_j \leftarrow P_j\_data['polygon']$
    \STATE $P_j\_name \leftarrow P_j\_data['name']$
    \STATE $P_j\_parent\_face \leftarrow P_j\_data['parent\_face']$
    \STATE $intersection \leftarrow P_i \cap P_j$
    \IF{$intersection.area > \epsilon$ AND $\neg intersection.is\_empty$}
        \STATE $intersection\_found \leftarrow \text{True}$
        \STATE Proceed to Step 3c (Enhanced Depth Analysis)
    \ENDIF
\ENDFOR
\end{algorithmic}
\end{algorithm}

\subsubsection{Step 3c: Enhanced Depth Analysis with Face Association}

This is the core enhancement of the algorithm, featuring multi-point sampling and face association.

\begin{algorithm}
\caption{Enhanced Depth Analysis}
\begin{algorithmic}[1]
\STATE $interior\_point \leftarrow \text{find\_interior\_point}(intersection)$
\STATE Create 3×3 grid sampling points within intersection region:
\STATE $sample\_points \leftarrow \emptyset$
\STATE $(x_{min}, y_{min}, x_{max}, y_{max}) \leftarrow intersection.bounds$
\FOR{$i = 0$ to $2$}
    \FOR{$j = 0$ to $2$}
        \STATE $x \leftarrow x_{min} + (i + 0.5) \times \frac{x_{max} - x_{min}}{3}$
        \STATE $y \leftarrow y_{min} + (j + 0.5) \times \frac{y_{max} - y_{min}}{3}$
        \STATE $test\_point \leftarrow \text{Point}(x, y)$
        \IF{$intersection.contains(test\_point)$}
            \STATE $sample\_points \leftarrow sample\_points \cup \{test\_point\}$
        \ENDIF
    \ENDFOR
\ENDFOR
\end{algorithmic}
\end{algorithm}

\begin{algorithm}
\caption{Farthest Point Detection}
\begin{algorithmic}[1]
\STATE $P_i\_max\_depth \leftarrow -\infty$
\STATE $P_j\_max\_depth \leftarrow -\infty$
\FOR{each $sample\_point$ in $sample\_points$}
    \STATE $P_i\_intersection\_3d \leftarrow \text{intersect\_line\_with\_face}(sample\_point, \vec{n}_{proj}, P_i\_parent\_face)$
    \STATE $P_j\_intersection\_3d \leftarrow \text{intersect\_line\_with\_face}(sample\_point, \vec{n}_{proj}, P_j\_parent\_face)$
    \STATE $P_i\_depth \leftarrow \text{calculate\_depth\_along\_normal}(P_i\_intersection\_3d, \vec{n}_{proj})$
    \STATE $P_j\_depth \leftarrow \text{calculate\_depth\_along\_normal}(P_j\_intersection\_3d, \vec{n}_{proj})$
    \IF{$P_i\_depth > P_i\_max\_depth$}
        \STATE $P_i\_max\_depth \leftarrow P_i\_depth$
    \ENDIF
    \IF{$P_j\_depth > P_j\_max\_depth$}
        \STATE $P_j\_max\_depth \leftarrow P_j\_depth$
    \ENDIF
\ENDFOR
\end{algorithmic}
\end{algorithm}

\textbf{Mathematical Foundation}: 
\begin{align}
\text{depth}(\vec{p}) &= \vec{p} \cdot \hat{n}_{proj} \\
\text{where } \hat{n}_{proj} &= \frac{\vec{n}_{proj}}{|\vec{n}_{proj}|}
\end{align}

\subsubsection{Step 3c.3: Face Association and Intersection Storage}

\begin{algorithm}
\caption{Face Association for Intersections}
\begin{algorithmic}[1]
\IF{$P_i\_max\_depth > P_j\_max\_depth$}
    \STATE $farthest\_face \leftarrow P_i\_parent\_face$
    \STATE $farthest\_name \leftarrow P_i\_name$
\ELSE
    \STATE $farthest\_face \leftarrow P_j\_parent\_face$
    \STATE $farthest\_name \leftarrow P_j\_name$
\ENDIF
\STATE $intersection\_name \leftarrow \text{"Intersection\_"} + P_i\_name + \text{"\_"} + P_j\_name$
\STATE $intersection\_data \leftarrow \{$
\STATE \quad $'polygon': intersection,$
\STATE \quad $'name': intersection\_name,$
\STATE \quad $'normal': \text{"intersection"},$
\STATE \quad $'parent\_face': farthest\_face,$
\STATE \quad $'associated\_face\_name': farthest\_name,$
\STATE \quad $'original\_index': -1$
\STATE $\}$
\STATE $\text{array\_C} \leftarrow \text{array\_C} \cup \{intersection\_data\}$
\end{algorithmic}
\end{algorithm}

\textbf{Key Enhancement}: Each intersection rectangle is immediately associated with the face that has the farthest point along the projection normal direction.

\subsubsection{Step 3c.4: Boolean Operations Based on Depth}

\begin{algorithm}
\caption{Depth-Based Boolean Operations}
\begin{algorithmic}[1]
\IF{$P_i\_max\_depth > P_j\_max\_depth$}
    \STATE // $P_i$ is farther: modify $P_j = P_j - P_i$
    \STATE $new\_P_j \leftarrow P_j \setminus P_i$
    \IF{$new\_P_j.area > \epsilon$}
        \STATE Update $P_j\_data$ in $\text{array\_B}$ with $new\_P_j$
    \ELSE
        \STATE Remove $P_j\_data$ from $\text{array\_B}$
    \ENDIF
\ELSE
    \STATE // $P_j$ is farther: modify $P_i = P_i - P_j$
    \STATE $new\_P_i \leftarrow P_i \setminus P_j$
    \IF{$new\_P_i.area > \epsilon$}
        \STATE $P_i \leftarrow new\_P_i$
        \STATE Update $P_i\_data$ with $new\_P_i$
    \ELSE
        \STATE $P_i$ is completely consumed
        \STATE \textbf{break} // Exit intersection testing loop
    \ENDIF
\ENDIF
\end{algorithmic}
\end{algorithm}

\subsection{Step 3d: Final Polygon Placement}

\begin{algorithm}
\caption{Final Polygon Placement}
\begin{algorithmic}[1]
\STATE // Note: Intersections were already added to array\_C in Step 3c.3
\IF{$intersection\_found = \text{True}$}
    \IF{$P_i\_data['polygon'].area > \epsilon$}
        \STATE $\text{array\_C} \leftarrow \text{array\_C} \cup \{P_i\_data\}$
        \STATE // Remaining $P_i$ goes to array\_C (classified)
    \ENDIF
    \STATE // If $P_i$ was completely consumed, it's not added anywhere
\ELSE
    \STATE $\text{array\_B} \leftarrow \text{array\_B} \cup \{P_i\_data\}$
    \STATE // No intersections: $P_i$ goes to array\_B for future comparisons
\ENDIF
\ENDWHILE
\end{algorithmic}
\end{algorithm}

\section{Algorithm Enhancements}

\subsection{Multi-Point Sampling Enhancement}

Traditional algorithms use a single interior point for depth comparison. This enhanced algorithm uses a 3×3 grid sampling approach:

\begin{equation}
\text{Sample Points} = \left\{(x_i, y_j) : i,j \in \{0,1,2\}\right\}
\end{equation}

where:
\begin{align}
x_i &= x_{min} + (i + 0.5) \times \frac{x_{max} - x_{min}}{3} \\
y_j &= y_{min} + (j + 0.5) \times \frac{y_{max} - y_{min}}{3}
\end{align}

\subsection{Farthest Point Detection}

Instead of using a single depth value, the algorithm finds the maximum depth across all sample points:

\begin{equation}
\text{max\_depth}(P) = \max_{s \in \text{sample\_points}} \left\{ \text{depth}(\text{intersect\_line\_with\_face}(s, \vec{n}_{proj}, P_{\text{parent}})) \right\}
\end{equation}

\subsection{Face Association Logic}

The algorithm associates each intersection with the face that contributes the farthest point:

\begin{equation}
\text{associated\_face} = \begin{cases}
P_i\_parent\_face & \text{if } P_i\_max\_depth > P_j\_max\_depth \\
P_j\_parent\_face & \text{if } P_j\_max\_depth > P_i\_max\_depth
\end{cases}
\end{equation}

\section{Expected Results}

\subsection{Intersection Rectangle Creation}

For a scenario where Face 11 intersects with both Face 4 and Face 6:

\begin{enumerate}
    \item \textbf{First intersection} (Face 11 ∩ Face 4):
    \begin{itemize}
        \item Algorithm determines which face (11 or 4) has the farthest point
        \item Creates \texttt{Intersection\_Face\_11\_Face\_4} associated with the farthest face
        \item Adds it to array\_C immediately
    \end{itemize}
    
    \item \textbf{Second intersection} (Face 11 ∩ Face 6):
    \begin{itemize}
        \item Algorithm determines which face (11 or 6) has the farthest point
        \item Creates \texttt{Intersection\_Face\_11\_Face\_6} associated with the farthest face
        \item Adds it to array\_C immediately
    \end{itemize}
\end{enumerate}

\subsection{Final Array Contents}

\begin{itemize}
    \item \textbf{array\_B}: Contains final polygons that didn't intersect with anything
    \item \textbf{array\_C}: Contains:
    \begin{itemize}
        \item All intersection rectangles with face associations
        \item Remaining classified polygon pieces
    \end{itemize}
\end{itemize}

\section{Algorithmic Complexity}

\subsection{Time Complexity}
\begin{itemize}
    \item \textbf{Main loop}: $O(n)$ where $n$ is the number of polygons
    \item \textbf{Intersection testing}: $O(n \times m)$ where $m$ is the average size of array\_B
    \item \textbf{Multi-point sampling}: $O(9 \times k)$ where $k$ is the cost of 3D line-face intersection
    \item \textbf{Overall}: $O(n^2 \times k)$ in worst case
\end{itemize}

\subsection{Space Complexity}
\begin{itemize}
    \item \textbf{Polygon storage}: $O(n)$ for all polygon data structures
    \item \textbf{Sample points}: $O(9)$ per intersection (constant)
    \item \textbf{Overall}: $O(n)$ space complexity
\end{itemize}

\section{Key Algorithmic Improvements}

\subsection{Problem Solved: Multiple Intersection Rectangles}

\textbf{Previous Issue}: When Face 11 intersected with both Face 4 and Face 6, only one intersection rectangle was created because the algorithm modified polygons before storing intersections.

\textbf{Solution}: The algorithm now adds intersection rectangles to array\_C \textbf{immediately} when found, \textbf{before} applying boolean operations:

\begin{lstlisting}[language=Python]
# This happens for EACH intersection found:
# Face 11 ∩ Face 4 → Creates "Intersection_Face_11_Face_4" in array_C
# Face 11 ∩ Face 6 → Creates "Intersection_Face_11_Face_6" in array_C
\end{lstlisting}

\subsection{Enhanced Depth Analysis}

The algorithm uses sophisticated 3D depth analysis with:
\begin{itemize}
    \item Multi-point sampling for accuracy
    \item Farthest point detection along projection normal
    \item Proper face association for traceability
\end{itemize}

\section{Conclusion}

This enhanced algorithm provides:
\begin{enumerate}
    \item \textbf{Accurate depth classification} using multi-point sampling
    \item \textbf{Separate intersection tracking} for proper visualization
    \item \textbf{Face association} for each intersection region
    \item \textbf{Robust 3D geometric analysis} with farthest point detection
\end{enumerate}

The result is a sophisticated polygon classification system that correctly handles complex 3D boolean operations with proper depth-based analysis and comprehensive intersection tracking.

\end{document}
